% Es werden nur die Abkürzungen aufgelistet, die mit \ac definiert und auch benutzt wurden. 
%
% \acro{VERSIS}{Versicherungsinformationssystem\acroextra{ (Bestandsführungssystem)}}
% Ergibt in der Liste: VERSIS Versicherungsinformationssystem (Bestandsführungssystem)
% Im Text aber: \ac{VERSIS} -> Versicherungsinformationssystem (VERSIS)

% Hinweis: allgemein bekannte Abkürzungen wie z.B. bzw. u.a. müssen nicht ins Abkürzungsverzeichnis aufgenommen werden
% Hinweis: allgemein bekannte IT-Begriffe wie Datenbank oder Programmiersprache müssen nicht erläutert werden,
%          aber ggfs. Fachbegriffe aus der Domäne des Prüflings (z.B. Versicherung)

\begin{acronym}	
	\acro{API}{Application Programming Interface}
	\acro{BTC}{Business Technology Consulting AG}
	\acro{CSS}{Cascading Style Sheets}
	\acro{HTML}{Hypertext Markup Language}
	\acro{HTML DOM}{\textit{HTML} Document Object Model}
	\acro{IDE}{Integrated Development Environment}
	\acro{Java EE}{Java Enterprise Edition}
	\acro{JSON}{JavaScript Object Notation}
	\acro{SQL}{Structured Query Language}
	\acro{TSO}{Transmission System Operator}
	\acro{TSO BSS}{\acs{TSO} Business Service Suit}
	\acro{UI}{User Interface}	
	\acro{VMS}{Vertragsmanagementsystem}
	\acro{XML}{Extensible Markup Language}
\end{acronym}
